%%%%%%%%%%%%%%%%%%%%%%%%%%%%%%%%%%%%%%%%%%%%%%%%%%%%%%%%%%%%%%%%%
%%%%%%%%%%  BiSO report starter - human-owned after creation  %%%
%%%%%%%%%%  Edit report text and commentary in this file.  %%%
%%%%%%%%%%%%%%%%%%%%%%%%%%%%%%%%%%%%%%%%%%%%%%%%%%%%%%%%%%%%%%%%%
\documentclass[french,11pt]{dibiso/biso}
\input{generated_variables.tex}
\ifdefined\tuttiBibliographyFile
  \usepackage{biblatex}
  \addbibresource{\tuttiBibliographyFile}
\fi
\title{Bilan \\ Science \\ Ouverte}
\subtitle{\textbf{Laboratoire \tuttiCfgEntityAcronym} \\
  \medskip
  \tuttiCfgEntityFullName}
\titledescription{\tuttiReportTitle}
\reporter{}
\ifdefined\tuttiCfgYear
  \date{Année \tuttiCfgYear}
\else
  \date{}
\fi
\providecommand{\makebiblio}{}
\renewcommand{\makebiblio}{%
\ifdefined\tuttiBibliographyFile
  \clearpage
  \nocite{*}
  {\footnotesize\printbibliography}
\fi
}
\begin{document}
\maketitle
\tableofcontents
\clearpage
\section{Introduction}
\subsection*{Méthodologie et objectifs}
Le BiSO est un bilan de science ouverte établi à partir des publications
présentes dans HAL. Sauf mention particulière, les données utilisées pour les
graphiques proviennent de la collection HAL \texttt{\tuttiCfgEntityId} et de
l’année \tuttiCfgYear.

% Ecrivez ici vos explications, objectifs et éléments de contexte.





% Fin du commentaire.

\subsection*{Calendrier annuel}
% Ecrivez ici les étapes et informations propres à votre laboratoire.





% Fin du commentaire.

\subsection*{Déposer dans une archive ouverte}
% Ecrivez ici vos recommandations et liens utiles.





% Fin du commentaire.

\clearpage
\input{generated_body.tex}
\input{generated_bibliography.tex}
\end{document}
